% =============================================================
% Section 5 -- Hardware validation of the model-necessity findings
% on the vessel-navigation task. IEEEtran (journal). Standalone
% compilable excerpt, companion to methodology.tex and
% model_necessity_study.tex (\S\ref{sec:model-necessity}): that
% section established the Jacobian-mismatch mechanism and its
% closed-loop consequences offline (Monte Carlo, contact flag never
% active on the accepted plan). This section reports the matching
% closed-loop hardware campaign (n=5 runs/condition, floor240
% vessel geometry) that was not available when \S\ref{sec:model-necessity}
% was written, and that does exercise the near-wall region directly.
% =============================================================
\documentclass[journal]{IEEEtran}

\usepackage{amsmath}
\usepackage{amssymb}
\usepackage{amsfonts}
\usepackage{booktabs}

\begin{document}

\section{Hardware Validation of the Model-Necessity Findings}
\label{sec:hw-validation}

\subsection{Conditions and Metrics}
\label{sec:hw-conditions}

Five conditions share the controller, task weights, and rate/state
limits of \S\ref{sec:control} and differ only in the Jacobian source
and, for one condition, the reference trajectory itself:
\begin{center}
\begin{tabular}{@{}lll@{}}
\toprule
Condition & Jacobian along horizon & Trajectory \\
\midrule
MPC-SJ & $J_c(z_k)$, scheduled every tick & $\mathcal{T}_c$ \\
$J_{nc}\!\to\!\mathcal{T}_c$ & $J_{nc}(z_k)$, scheduled every tick & $\mathcal{T}_c$ \\
MPC-FJ & $J_c(z_0)$, frozen at $k{=}0$ & $\mathcal{T}_c$ \\
$J_{nc}\!\to\!\mathcal{T}_{nc}$ & $J_{nc}(z_k)$, scheduled every tick & $\mathcal{T}_{nc}$ \\
INV-7 & $J_c(z_k)$, scheduled every tick & $\mathcal{T}_c$ \\
\bottomrule
\end{tabular}
\end{center}
$\mathcal{T}_c$ is the reference trajectory of \S\ref{sec:mn-design},
designed against the contact-aware beam model; $\mathcal{T}_{nc}$ is
the corresponding trajectory designed against the no-contact model.
Every condition is run $5\times$ to \texttt{path\_complete} or to a
hard tracking-error abort at $7\,\mathrm{mm}$; RMSE and lag are
reported per \eqref{eq:e-rms} over three fixed, physical-progress
windows: the whole run, an early low-authority region
$s\in[12,17]\,\mathrm{mm}$ (``region A''), and the near-wall region
$s\in[55,65]\,\mathrm{mm}$ (``region C'') identified in
\S\ref{sec:model-necessity} as where $J_c$ and $J_{nc}$ diverge most.
Run-level mean $\pm$ 95\% CI uses $t_{0.975,4}\approx2.776$ on the 5
run-level values (\S\ref{sec:metrics}); where a condition does not
reach region C in all 5 runs, $n$ is reported explicitly and no
interval is computed.

\subsection{Contact-Aware vs.\ No-Contact Jacobian}
\label{sec:hw-exp1}

\begin{center}
\begin{tabular}{@{}lccc@{}}
\toprule
& whole-run $E_{\mathrm{RMS}}$ & region A & region C \\
\midrule
MPC-SJ ($J_c$) & $0.962\pm0.057$ & $0.309\pm0.100$ & $0.619\pm0.073$ \\
$J_{nc}\!\to\!\mathcal{T}_c$ & $0.996\pm0.054$ & $0.288\pm0.071$ & $\mathbf{0.367\pm0.060}$ \\
\bottomrule
\end{tabular}
\end{center}
(mm; $n=5$ throughout.) $J_c$ is more accurate whole-run (non-overlapping
95\% CIs) and region A is statistically indistinguishable between the
two. Region C reverses this: $J_{nc}$ is \emph{more} accurate than
$J_c$ there, by a wide, non-overlapping margin ($0.367\pm0.060$ vs.\
$0.619\pm0.073\,\mathrm{mm}$) --- the opposite of the whole-run result
and the opposite of what $E_J(s)$ (\S\ref{sec:mn-exp1}) would predict
from local model fidelity alone, since $J_{nc}$ is the strictly worse
local approximation to the plant everywhere on $\mathcal{T}_c$,
including region C. Mean lag magnitude favors $J_c$ in both regimes
($0.169\pm0.035$ vs.\ $0.187\pm0.032\,\mathrm{mm}$ whole-run;
$0.532\pm0.065$ vs.\ $0.736\pm0.034\,\mathrm{mm}$ in region C), so the
reversal is specific to tracking accuracy, not a general degradation
of $J_c$'s closed-loop behavior near the wall.

\subsection{Scheduled vs.\ Frozen Jacobian}
\label{sec:hw-exp2}

\begin{center}
\begin{tabular}{@{}lcc@{}}
\toprule
& whole-run $E_{\mathrm{RMS}}$ & region A \\
\midrule
MPC-SJ & $0.962\pm0.057$ & $0.309\pm0.100$ \\
MPC-FJ & $2.371\pm0.133$ & $1.456\pm0.244$ \\
\bottomrule
\end{tabular}
\end{center}
(mm; $n=5$.) MPC-FJ is categorically worse, by non-overlapping margins,
at every progress window tested, and reaches region C in \emph{none}
of its 5 runs (all abort on tracking error at $s\approx31$--$41\,\mathrm{mm}$).
The degradation is already visible in region A, only
$12$--$17\,\mathrm{mm}$ into the path: freezing $J$ at $k{=}0$ costs
accuracy well before the beam reaches the contact-affected region
identified in \S\ref{sec:mn-exp1}, confirming that re-evaluating $J$
along the path matters independently of, and prior to, any
contact-specific effect.

\subsection{Contact-Modeled vs.\ No-Contact-Modeled Reference Trajectory}
\label{sec:hw-exp3}

\begin{center}
\begin{tabular}{@{}lccc@{}}
\toprule
& whole-run $E_{\mathrm{RMS}}$ & region A & region C \\
\midrule
$\mathcal{T}_c$ (MPC-SJ) & $0.962\pm0.057$ & $0.309\pm0.100$ & $0.619\pm0.073$ \\
$\mathcal{T}_{nc}$ & $2.083\pm0.055$ & $0.281\pm0.019$ & $0.516\pm0.031$ \\
\bottomrule
\end{tabular}
\end{center}
(mm; $n=5$.) Designing the reference trajectory against the
contact-aware model is markedly more accurate whole-run
(non-overlapping CIs, more than $2\times$), with region A again
statistically indistinguishable. Region C again reverses
($0.516\pm0.031$ vs.\ $0.619\pm0.073\,\mathrm{mm}$, non-overlapping),
but $\mathcal{T}_{nc}$ remains clearly less accurate there than
$J_{nc}\!\to\!\mathcal{T}_c$ from \S\ref{sec:hw-exp1}
($0.516\pm0.031$ vs.\ $0.367\pm0.060\,\mathrm{mm}$): using the
no-contact Jacobian is only beneficial near the wall when applied to
the \emph{contact-designed} trajectory, not to a trajectory already
shaped to avoid the contact-model's predicted wall interaction. The
near-wall reversal is therefore a property of the $(J,\mathcal{T})$
pairing, not of omitting contact modeling in general.

\subsection{MPC vs.\ Matched-Authority Inverse-Jacobian Control (INV-7)}
\label{sec:hw-exp4}

\begin{center}
\begin{tabular}{@{}lcc@{}}
\toprule
& whole-run $E_{\mathrm{RMS}}$ & region A \\
\midrule
MPC-SJ & $0.962\pm0.057$ & $0.309\pm0.100\;(n{=}5)$ \\
INV-7 & $3.037\pm2.217$ & $0.420\pm0.009\;(n{=}3)$ \\
\bottomrule
\end{tabular}
\end{center}
(mm.) MPC-SJ is more accurate whole-run; INV-7's wide interval reflects
2 of 5 runs terminating within $44$ ticks on an unconstrained-nullspace
drift into a workspace or exclusion bound (\S\ref{sec:inv7}), not
measurement noise --- the instability is itself part of the
comparison and is not averaged away by excluding those runs. Of
INV-7's remaining 3 runs, only 1 reaches region C at all, versus 5 of
5 for MPC-SJ; no region-C interval is reported for INV-7 for that
reason. MPC's hard in-QP constraint enforcement (\S\ref{sec:mpc})
keeps it running long enough to reach the contact-affected region on
every trial, which INV-7's soft, preview-based bound handling
(\S\ref{sec:inv7}) does not reliably do.

\subsection{Synthesis}
\label{sec:hw-synthesis}

Three of the four comparisons above are unqualified: scheduling $J$
beats freezing it, a contact-designed trajectory beats a no-contact-designed
one, and MPC beats INV-7, all without exception across progress
windows. The fourth is not unqualified, and should not be reported as
one: the contact-aware Jacobian is more accurate everywhere tested
\emph{except} the near-wall region, where the no-contact Jacobian wins
by the same margin the contact Jacobian wins everywhere else. This is
the single result in this section that contradicts a naive reading of
\S\ref{sec:mn-exp1}'s local Jacobian-mismatch characterization --- the
locally worse model produces the better closed-loop outcome exactly
where the mismatch is largest. \S\ref{sec:hw-exp3}'s result that this
only holds when the no-contact Jacobian is paired with the
contact-designed trajectory narrows the explanation to a property of
that specific pairing, not a general claim that contact modeling is
unnecessary near the wall.

\end{document}
